% ══════════════════════════════════════════════════════════════════════════════
%  EXPERIMENT 1 — MATTER HAS MASS AND TAKES UP SPACE                   (Lesson 3)
%  Content from the lesson booklet's Lesson 3, plus the dry ice demonstration
%  from the unit plan.
% ══════════════════════════════════════════════════════════════════════════════
\experimentdiv{Matter Has Mass and Takes Up Space}{Dry ice demonstration \textbullet{} Does air have mass? \textbullet{} Does air take up space?}{3}\label{sec:1}
\lessoncodes{VC2S8U05 \textbullet{} VC2S8I01, I02, I04, I06}

\vocabhead
\begin{ctab}
\begin{tabular}{|SW{4.2cm}|W{11.6cm}|}
  \hline\hd & \thead{Word} & \thead{Definition (copy from the slide)} \\ \hline
  & \scaf{\textbf{Matter}} & \scaf{Anything that has mass and takes up space.} \\ \hline
  & \textbf{Observation} & \\ \hline
  & \textbf{Inference} & \\ \hline
  & \textbf{Sublimation} & \\ \hline
  & \textbf{Tare} & \\ \hline
\end{tabular}
\end{ctab}

\begin{lipanel}
\inquiryli{Use a balance to take precise readings, and record them in a table.}
\understandli{Explain, with evidence, that gases are matter: they have mass and take up space.}
\sctitle
\begin{itemize}[leftmargin=1.4em, itemsep=0pt, topsep=1pt]
  \item I can tell an observation from an inference.
  \item I can measure a small change in mass and calculate the difference and the average.
  \item I can use my results to explain that air has mass and takes up space.
\end{itemize}
\end{lipanel}

\begin{skillsfocus}
\focus{QP1}{Predict what will happen to the ball's mass, and why.}
\focus{PC4}{Use the same balance, tared, and keep every decimal place.}
\focus{PC6}{Record observations of the demonstration and both tests.}
\focus{PA1}{Record masses in a table with units in the headings.}
\focus{PA3}{Calculate each increase in mass, and the average increase.}
\focus{EV1}{Say whether your results support your prediction.}
\focus{EV3}{Name what could make the mass readings less accurate.}
\end{skillsfocus}

% ─── TECHNIQUE ────────────────────────────────────────────────────────────────
\exppart{Technique spotlight}{Blue}

\begin{technique}{Observation or inference?}{PC6}
An \textbf{observation} is what you can see, hear, feel or measure: \emph{the balloon got bigger}. An \textbf{inference} is an explanation of an observation: \emph{gas came off the dry ice and filled the balloon}. Record observations first. Keep inferences for the analysis, and back each one with an observation.
\end{technique}

\begin{workedex}{Sorting statements}
A teacher drops dry ice into warm water. Sort the statements.

\smallskip
\begin{tabular}{@{}p{9.4cm}l@{}}
\emph{White fog spills over the side of the beaker.} & \textbf{Observation} --- it can be seen. \\
\emph{The water is bubbling.} & \textbf{Observation} --- it can be seen and heard. \\
\emph{The bubbles are carbon dioxide gas.} & \textbf{Inference} --- it explains what is seen. \\
\end{tabular}
\end{workedex}

\begin{wedo}{Sorting statements}
Write \emph{observation} or \emph{inference} for each statement.
\begin{ctab}
\begin{tabular}{|SW{10.4cm}|W{4.6cm}|}
  \hline\hd & \thead{Statement} & \thead{Observation or inference?} \\ \hline
  & The piece of dry ice is getting smaller. & \\ \hline
  & The dry ice is turning straight into a gas. & \\ \hline
  & The balloon on the flask inflated. & \\ \hline
\end{tabular}
\end{ctab}
\end{wedo}

\begin{technique}{Measuring a small change in mass}{PC4}
Pumping air into a ball changes its mass by only a few grams, so every reading has to be as precise as possible.
\begin{enumerate}[leftmargin=1.6em, itemsep=1pt, topsep=3pt]
  \item Use the \textbf{same} balance for the before and after readings.
  \item Check it reads 0.00\,g with nothing on it, and keep every decimal place.
  \item Find the change \textbf{by difference}: mass after $-$ mass before.
  \item Repeat the test three times and average the changes.
\end{enumerate}
\end{technique}

\begin{workedex}{Mass by difference}
A student forgets to tare the balance. An empty beaker reads 52.40\,g. With salt in it, it reads 67.15\,g. What is the mass of the salt?

\medskip
\textbf{Mass of salt} = mass with salt $-$ mass of empty beaker $= 67.15 - 52.40 = \mathbf{14.75\ g}$
\end{workedex}

\begin{wedo}{Mass by difference}
An empty plastic cup has a mass of 8.30\,g. Filled with sand, it has a mass of 126.80\,g. What is the mass of the sand?
\workbox{3cm}
\end{wedo}

\begin{youdobox}[Changes in mass]
\taskpair{A ball had a mass of 410.50\,g before pumping and 413.10\,g after. Find the mass of air added.}{3.0cm}
         {Three tests added 2.60\,g, 2.75\,g and 2.65\,g of air. Find the average mass of air added.}{3.0cm}
\end{youdobox}

% ─── PLAN AND SAFETY / METHOD ────────────────────────────────────────────────
\exppart{Plan, safety and method}{Teal}

\textbf{Aim:} To show that air has mass and takes up space.

\needspace{10\baselineskip}
\task Will a ball weigh more, less or the same after you pump more air into it? Write your prediction and explain it. \skr{QP1}
\answerbox{5cm}

\subsubsection*{Equipment}
\begin{multicols}{2}
\begin{itemize}[leftmargin=1.2em, itemsep=0pt, topsep=0pt]
  \item soft ball with a valve (basketball or netball)
  \item ball pump with a needle
  \item electronic balance (reads to 0.01\,g)
  \item large clear tub of water
  \item clear plastic cup
  \item paper towel and sticky tape
\end{itemize}
\end{multicols}

\begin{infobox}{Safety}
\textbf{Demonstration:} only your teacher handles dry ice. At $-78$\,\textdegree C it burns skin on contact, and it is never sealed in a closed container. Watch from where your teacher tells you to stand.\par
\textbf{Your tests:} wipe up spilt water straight away --- a wet floor is a slip hazard. Keep the pump needle pointed away from people, and keep water away from the balance and its power cord.
\end{infobox}

\subsubsection*{Demonstration --- Dry ice (your teacher)}
Your teacher puts a piece of dry ice (solid carbon dioxide) on the bench, drops some into warm water, and puts some in a flask with a balloon stretched over its neck. Record what you observe in Table 1.1.

\subsubsection*{Method --- Part 1: Does air have mass?}
\begin{enumerate}[itemsep=2pt]
  \item Let some air out of the ball so it is soft but still round.
  \item Check the balance reads 0.00\,g, put the ball on it, and record its mass.
  \item Pump the ball until it is firm. Count the number of pumps.
  \item Weigh the ball again and record its mass.
  \item Repeat the whole test two more times.
\end{enumerate}

\subsubsection*{Method --- Part 2: Does air take up space?}
\begin{enumerate}[itemsep=2pt]
  \item Scrunch a paper towel and tape it into the bottom of the cup so it cannot fall out.
  \item Turn the cup upside down and push it straight down to the bottom of the tub. Do not tilt it.
  \item Lift it straight out, and feel the paper towel.
  \item Now repeat, but tilt the cup while it is under water. Watch what comes out.
\end{enumerate}

% ─── RESULTS ─────────────────────────────────────────────────────────────────
\exppart{Results}{Amber}

\subsubsection*{Table 1.1 --- Dry ice demonstration \skr{PC6}}
\begin{ctab}
\begin{tabular}{|SW{3.4cm}|W{5.8cm}|W{5.8cm}|}
  \hline\hd & \thead{Dry ice \dots} & \thead{What I observed} & \thead{What I infer from it} \\ \hline
  & \scaf{e.g. held over a candle flame} & \scaf{The flame went out.} & \scaf{The gas does not let things burn.} \\ \hline
  & on the bench & & \\ \hline
  & in warm water & & \\ \hline
  & in a flask with a balloon & & \\ \hline
\end{tabular}
\end{ctab}

\subsubsection*{Table 1.2 --- Does air have mass? \skr{PA1}}
\begin{ctab}
\begin{tabular}{|SZ{1.6cm}|Z{3.3cm}|Z{2.0cm}|Z{3.3cm}|Z{4.4cm}|}
  \hline\hd & \thead{Test} & \thead{Mass before pumping (g)} & \thead{Pumps} & \thead{Mass after pumping (g)} & \thead{Increase in mass (g)} \\ \hline
  & \scaf{e.g.} & \scaf{598.20} & \scaf{20} & \scaf{601.65} & \scaf{$601.65 - 598.20 = 3.45$} \\ \hline
  & 1 & & & & \\ \hline
  & 2 & & & & \\ \hline
  & 3 & & & & \\ \hline
  \multicolumn{5}{|r|}{\rule{0pt}{1.5cm}\headingfont\bfseries Average increase in mass (g)} & \\ \hline
\end{tabular}
\end{ctab}

\subsubsection*{Table 1.3 --- Does air take up space?}
\begin{ctab}
\begin{tabular}{|SW{5.0cm}|W{10.0cm}|}
  \hline\hd & \thead{Test} & \thead{What I observed} \\ \hline
  & \scaf{e.g. an empty cup pushed down} & \scaf{The water only came part of the way up inside the cup.} \\ \hline
  & Cup pushed straight down: the paper towel & \\ \hline
  & Cup tilted under water & \\ \hline
\end{tabular}
\end{ctab}

% ─── ANALYSIS ────────────────────────────────────────────────────────────────
\exppart{Analysis and evaluation}{Coral}

\task Did the ball's mass go up, go down, or stay the same when air was pumped in? Was your prediction correct? \skr{EV1}
\answerbox{5cm}

\needspace{10\baselineskip}
\task Explain what your results from Table 1.2 show about air. Use the word \emph{mass}.
\answerbox{5cm}

\needspace{10\baselineskip}
\task Explain why the paper towel stayed dry. Use the words \emph{air} and \emph{space}.
\answerbox{5cm}

\needspace{10\baselineskip}
\task What came out of the cup when you tilted it under water, and what does this show?
\answerbox{5cm}

\needspace{10\baselineskip}
\task The balloon on the flask inflated even though nothing was blown into it. Use the dry ice demonstration to explain where the gas came from, and what this shows about gases.
\answerbox{5cm}

\needspace{14\baselineskip}
\task Evaluate the method for Part 1. \skr{EV3}
\pt{a}Why did you repeat the ball test three times?
\answerbox{5cm}
\pt{b}Identify one thing that could have made your mass readings less accurate.
\answerbox{5cm}
\pt{c}Suggest one improvement to Part 1.
\answerbox{5cm}

\begin{yr8box}
\begin{keepext}
\exttask One litre of air has a mass of about 1.2\,g. A ball gained 3.6\,g when it was pumped up.
\pt{a}About how many litres of extra air were pushed into the ball?
\workbox{3cm}
\pt{b}The ball itself holds only about 7 litres. Explain how it can hold that much extra air.
\answerbox{5cm}
\end{keepext}
\begin{keepext}
\exttask A classroom is 8\,m long, 7\,m wide and 3\,m high. One cubic metre of air has a mass of about 1.2\,kg. Calculate the volume of the room and the mass of air in it.
\workbox{4cm}
\end{keepext}
\begin{keepext}
\exttask A student weighs a balloon before and after blowing it up and finds almost no change. She concludes that air has no mass. Suggest why her result might have been misleading, and how Part 1 of your experiment avoided the problem.
\answerbox{5cm}
\end{keepext}
\end{yr8box}

% ─── SKILL CHECK ─────────────────────────────────────────────────────────────
\exppart{Skill check}{NavyBlue}

\begin{skillcheck}
\scrow{QP1}{I can write a prediction and explain it.}
\scrow{PC4}{I used the same balance, checked it read zero, and kept every decimal place.}
\scrow{PC6}{I can record observations, and tell them apart from inferences.}
\scrow{PA1}{I recorded masses in a table with units in the headings.}
\scrow{PA3}{I can calculate a change in mass by difference, and an average.}
\scrow{EV1}{I can say whether my results supported my prediction.}
\scrow{EV3}{I can name what could make a reading less accurate, and suggest an improvement.}
\end{skillcheck}

\task Which skill will you focus on in Experiment 2? Explain what you will do differently.
\answerbox{5cm}
